\chapter{Formato de los ficheros de ensayo}
\label{anx:csv}

Cada ensayo se graba en \fichero{cycles/AAAAMMDD-HHMMSS\_ciclo.csv}, codificado en UTF-8. Las líneas
que empiezan por \codigo{\#} forman la cabecera, con las constantes en pares \codigo{clave=valor};
después viene una línea de nombres de columna y una fila por muestra.

\begin{lstlisting}[numbers=none,caption={Comienzo de un fichero de la campaña del 21 de septiembre.}]
# ensayo 2026-09-21 18:08:46
# lc1_cuentas_por_N=14.41  lc1_cero_N=51.78  lc1_int16=0
# repeticiones_pedidas=1  B=1.000 rev (posicion)
# lc1_base ya lleva restada la tara del firmware; lc1_raw es cruda
t_ms,t_rel_s,rep,fase,pos_rev,lc1_raw,lc1_base,rpm,par_x10,ref_rpm,lim_a,lim_b,estado,work_us
\end{lstlisting}

\begin{longtable}{@{}llL{0.58\textwidth}@{}}
\caption{Columnas del fichero de ensayo.}\label{tab:csv}\\
\toprule
Columna & Unidad & Descripción \\
\midrule
\endfirsthead
\toprule
Columna & Unidad & Descripción \\
\midrule
\endhead
t\_ms & ms & Reloj del microcontrolador \\
t\_rel\_s & s & Tiempo desde la primera muestra del ensayo \\
rep & -- & Repetición en curso (desde 1) \\
fase & -- & \codigo{hacia\_A} (subida) o \codigo{volviendo\_a\_B} (bajada). En los ensayos de la campaña preliminar, \codigo{hacia\_B} y \codigo{volviendo\_a\_A} \\
pos\_rev & rev & Posición del eje del motor integrada por el servidor \\
lc1\_raw & cuentas & Lectura bruta del ZSC31014 (una por lote) \\
lc1\_base & cuentas & Media móvil de 8 lecturas menos la tara del firmware \\
rpm & rpm & Velocidad medida por el accionamiento \\
par\_x10 & 0,1\,\ua & Par medido por el accionamiento \\
ref\_rpm & rpm & Consigna de velocidad (negativa hacia A, positiva hacia B) \\
lim\_a, lim\_b & 0/1 & Estado de los finales de carrera \\
estado & -- & Estado de la máquina (tabla~\ref{tab:estados}) \\
work\_us & µs & Tiempo de trabajo de la vuelta del lazo \\
\bottomrule
\end{longtable}

\section{Conversión y análisis}

La conversión a fuerza es $F\,[\si{\newton}] = \text{lc1\_base}/k + F_0$, con $k$ =
\codigo{lc1\_cuentas\_por\_N} y $F_0$ = \codigo{lc1\_cero\_N}. Conviene recordar que esa constante
es provisional (sección~\ref{sec:res-limitaciones}).

Para el análisis, el \emph{avance del ciclo} se obtiene restando la posición de la primera muestra,
que es el punto B donde arranca el ensayo, y corrigiendo el signo con el de la muestra más alejada:
\begin{equation}
  p_i = \big(\text{pos\_rev}_i - \text{pos\_rev}_0\big)\cdot\operatorname{signo}
        \big(\text{pos\_rev}_{\text{lejana}} - \text{pos\_rev}_0\big).
\end{equation}
Así el avance es siempre positivo hacia A, y los ensayos de las dos campañas se pueden comparar pese
a tener sentidos de recorrido opuestos.

\section{Notas sobre las cabeceras}

\begin{itemize}
  \item Los ensayos de la campaña preliminar se grabaron con una versión del servidor que todavía no
        escribía la calibración: en ellos aparece \codigo{lc1\_cuentas\_por\_N=0}.
  \item El texto ``(posicion)'' de la tercera línea se escribe siempre que el recorrido pedido sea
        mayor que cero, aunque el ensayo se haya hecho contra los finales de carrera, que es el caso
        de todos los ensayos de esta memoria. Es un defecto conocido de la cabecera y no afecta a los
        datos.
\end{itemize}
